\section*{Capítulo \ref{ch:g8:ohms-law} --- La resistencia y la ley de Ohm}

\begin{solution}{exo:g8:ohms-law:1}
Con cuánta \omterm{def:g2:pushes-and-pulls:force}{fuerza} se opone un componente a la corriente;
\omterm{def:g6:measurement-in-science:measuring}{unidad} el \omterm{def:g8:ohms-law:resistance}{ohmio}, símbolo $\Omega$. $U = R I$; \ $I
= U/R$; \ $R = U/I$.
\end{solution}

\begin{solution}{exo:g8:ohms-law:2}
$I = 6 \div 30 = \qty{0.2}{A}$; \ $U = 30 \times 0.5 = \qty{15}{V}$.
\end{solution}

\begin{solution}{exo:g8:ohms-law:3}
$R = 4.5 \div 0.09 = 50\,\Omega$.
\end{solution}

\begin{solution}{exo:g8:ohms-law:4}
Una recta que pasa por el origen --- proporcionalidad. Una recta más
empinada significa más \omterm{def:g8:voltage-voltmeter:voltage}{voltios} necesarios por
\omterm{def:g8:intensity-ammeter:intensity}{amperio}: mayor \omterm{def:g8:ohms-law:resistance}{resistencia}.
\end{solution}

\begin{solution}{exo:g8:ohms-law:5}
$3.0 \div 0.20 = 15$; $6.0 \div 0.40 = 15$. Que el cociente sea el
mismo en todas las filas \emph{es} la ley: que una sola constante $R$
sirva para todas las lecturas es lo que convierte a ``la
\omterm{def:g8:ohms-law:resistance}{resistencia}'' en una propiedad del componente.
\end{solution}

\begin{solution}{exo:g8:ohms-law:6}
La \omterm{def:g8:ohms-law:resistance}{resistencia} más pequeña bebe más: $230 \div 25 =
\qty{9.2}{A}$ frente a $230 \div 35 \approx \qty{6.6}{A}$.
\end{solution}

\begin{solution}{exo:g8:ohms-law:7}
Las lámparas \omterm{def:g4:series-circuits:series}{en serie} comparten una sola $I$; la \omterm{def:g8:voltage-voltmeter:voltage}{tensión} de cada
lámpara es $U = R I$ con la misma $I$ --- así que la $R$ mayor
multiplica hasta dar la $U$ mayor: el que más se opone se lleva la
parte más grande.
\end{solution}

\begin{solution}{exo:g8:ohms-law:8}
\omterm{def:g7:short-circuits-safety:device}{Cortocircuito}: $R$ casi nula, así que $I = U/R$ explota --- la
estampida es una división por casi nada. \omterm{ex:g3:first-electric-circuit:switch}{Interruptor} abierto:
$R$ más allá de toda medida, así que $I = U/R$ se derrumba hasta
cero --- el bloqueo total.
\end{solution}

\begin{solution}{exo:g8:ohms-law:9}
El \omterm{def:g8:ohms-law:resistor}{resistor} tiene que llevarse $9 - 2 = \qty{7}{V}$ a
\qty{0.02}{A}: $R = 7 \div 0.02 = 350\,\Omega$.
\end{solution}

\begin{solution}{exo:g8:ohms-law:10}
Primer punto: $R = 1 \div 0.25 = 4\,\Omega$; segundo: $4 \div 0.5 =
8\,\Omega$ --- ningún valor único encaja. La historia: la corriente que
crece calienta el filamento, y el metal caliente resiste más;
el retrato se curva porque el componente cambia mientras trabaja.
\end{solution}

\begin{solution}{exo:g8:ohms-law:11}
Un lazo, una $I$: los dos saltos suman el empujón de la
\omterm{def:g3:first-electric-circuit:battery}{pila}, $6 = 10 I + 20 I = 30 I$, así que $I = \qty{0.2}{A}$.
Entonces $U_{10} = 10 \times 0.2 = \qty{2}{V}$ y $U_{20} = 20 \times
0.2 = \qty{4}{V}$ --- que suman \qty{6}{V}, como exige la ley del
lazo.
\end{solution}

\begin{solution}{exo:g8:ohms-law:12}
Cada \omterm{def:g5:series-parallel-circuits:parallel}{rama} abarca \qty{6}{V}: $I_{10} = 6 \div 10 =
\qty{0.6}{A}$, $I_{20} = 6 \div 20 = \qty{0.3}{A}$; total
$\qty{0.9}{A}$. La sustituta: $R = 6 \div 0.9 \approx 6.7\,\Omega$ ---
\emph{menos} que cualquiera de los dos miembros. En lenguaje de
carreteras: dos carreteras entre los mismos pueblos llevan más tráfico
que cualquiera de ellas sola; cada carretera \omterm{def:g5:series-parallel-circuits:parallel}{en paralelo} que se añade
facilita el paso, así que la oposición del equipo cae por debajo de la
de su miembro más flojo.
\end{solution}

\begin{solution}{pb:g8:ohms-law:1}
\textbf{1.} La fuente, el \omterm{def:g8:intensity-ammeter:ammeter}{amperímetro} y el componente en un solo
lazo \omterm{def:g4:series-circuits:series}{en serie} --- el peaje en la carretera; el
\omterm{def:g8:voltage-voltmeter:voltmeter}{voltímetro} cruzado sobre el componente --- el topógrafo, a un
lado.
\textbf{2.} $R = 4.5 \div 0.15 = 30\,\Omega$.
\textbf{3.} El cinco por ciento de $33$ es aproximadamente $1.65$:
aceptable de $31.35$ a $34.65\,\Omega$. Los $30\,\Omega$ de la muestra
uno se quedan fuera: rechazada.
\textbf{4.} $I = 4.5 \div 150 = \qty{0.03}{A} = \qty{30.0}{mA}$.
\textbf{5.} Todas las filas devuelven $R = U/I = 30\,\Omega$ ($1.5 \div
0.05$, $3.0 \div 0.10$, \dots): cociente constante, retrato en línea
recta --- un solo valor sirve para todo: $30\,\Omega$.
\textbf{6.} Los cocientes van $5$, $6.7$, $8.3$, $10\,\Omega$: suben en
cada paso --- \omterm{def:g8:ohms-law:resistance}{resistencia} que crece a medida que crece la
corriente; no existe ninguna $R$ única.
\textbf{7.} A medida que crece la corriente, el filamento se
calienta hasta ponerse al rojo, y el metal caliente resiste más
--- el retrato se curva hacia arriba. La prueba de la recta del banco
admite entonces exactamente a los componentes cuya $R$ se queda quieta:
es, por construcción, un detector de \omterm{def:g8:ohms-law:resistor}{resistores}.
\textbf{8.} La cláusula de la \omterm{def:g3:temperature-thermometers:temperature}{temperatura} estable: la ley de Ohm
promete proporcionalidad \emph{a \omterm{def:g3:temperature-thermometers:temperature}{temperatura} estable} ---
calienta un componente a mitad del retrato y hasta un
\omterm{def:g8:ohms-law:resistor}{resistor} honrado se desvía.
\textbf{9.} $R = 12 \div 0.03 = 400\,\Omega$.
\textbf{10.} $U = R I = 25 \times 9.2 = \qty{230}{V}$ --- exactamente
la red: la \omterm{def:g8:ohms-law:resistance}{resistencia} absorbe precisamente lo que la ley de Ohm le
asigna, y el margen del \omterm{def:g7:short-circuits-safety:fuse}{fusible} está funcionando como se diseñó.
No hay ningún ``solo''; todo es aritmética.
\textbf{11.} El puente ha nacido para ser un cable de conexión --- la
carretera llana, que no gasta empujón. Mal usado entre los \omterm{def:g3:first-electric-circuit:battery}{bornes} de
una fuente, sus \omterm{def:g8:ohms-law:resistance}{ohmios} casi nulos lo convierten en el
\omterm{def:g7:short-circuits-safety:device}{cortocircuito} perfecto: el villano siempre fue, simplemente, un
cable muy bueno en el sitio equivocado.
\textbf{12.} Por ejemplo: ``La ley: $U = R I$, empujón igual a
oposición por desfile. Su retrato: una recta que pasa por el origen,
más empinada cuanto más se opone el componente. Y solo puede firmarla
un componente a \omterm{def:g3:temperature-thermometers:temperature}{temperatura} estable --- el \omterm{def:g5:heat-and-insulation:heat}{calor}
reescribe la \omterm{def:g8:ohms-law:resistance}{resistencia}, y los retratos tomados con fiebre no
demuestran nada''.
\end{solution}
